\section*{Apéndice: Índice de citas clave}

Lo que más vale la pena volver a ver de la entrevista no suele ser algún punto noticioso, sino estos juicios que Xie Saining (谢赛宁) comprimió una y otra vez en una sola frase. Su fuerza no viene de que las palabras sean elegantes, sino de que detrás de cada una hay una conciencia de largo plazo sobre el problema. A continuación se organiza como índice el conjunto de citas más importante de toda la entrevista junto con su contexto, para ubicar con rapidez su posición de pensamiento en revisiones posteriores.

\begin{longtable}{>{\raggedright\arraybackslash}p{0.34\textwidth} >{\raggedright\arraybackslash}p{0.58\textwidth}}
\toprule
Cita clave & Conciencia del problema correspondiente \\
\midrule
\endfirsthead
\toprule
Cita clave & Conciencia del problema correspondiente \\
\midrule
\endhead
\textit{``Los ojos son la única parte del cerebro expuesta al mundo real... resolver la visión no consiste en resolver la visión misma, sino en resolver la inteligencia misma''} & Esta es la premisa teórica más central de toda la entrevista. Eleva la visión de una tarea concreta a la interfaz entre la inteligencia y el mundo real, y también explica por qué Xie Saining terminó orientándose hacia el world model. \\
\textit{``La research nunca es un desarrollo lineal; una research que se desarrolla de forma lineal nunca es una buena research''} & Esta frase define su entendimiento básico del trabajo creativo. Una buena investigación debe incluir un proceso de desviación, conflicto y reformulación del problema; de lo contrario, eso indica que el problema mismo no es suficientemente profundo. \\
\textit{``La idea que se te ocurre al principio no es tu idea; la idea que surge durante la exploración es la que verdaderamente te pertenece''} & Subraya que la intuición que verdaderamente pertenece al investigador solo puede surgir en el proceso de hacer las cosas, no de pensarlo todo de una vez en la etapa de especulación. La exploración no es un acto preparatorio, sino la generación misma de la idea. \\
\textit{``Toda forma es ilusoria; quien ve que las formas no son formas, ve al Tathagata''} & En el contexto de la entrevista, esto no es una exhibición de una cita religiosa, sino una descripción del research taste: no dejarse engañar por la apariencia de los artículos, el empaque terminológico y los indicadores parciales, sino indagar las variables sustanciales que hay detrás. \\
\textit{``Los LLM terminarán marchitándose... los viejos soldados no mueren, solo se desvanecen''} & El punto central de este juicio no es negar los modelos de lenguaje, sino negar que se tome a los LLM como el punto final. Apunta a las limitaciones estructurales de los modelos de lenguaje en la representación del mundo real y en la capacidad de planeación. \\
\textit{``El pasado fue la era de download internet; el presente es la era de download human''} & Redefine la base de datos de la siguiente etapa de la IA. El internet de texto ya no es el único centro; el video, la percepción, la conducta y las trayectorias de los seres humanos en el mundo se convertirán en fuentes de entrenamiento más determinantes. \\
\textit{``Poder volver a crear una ardilla es mucho más grandioso que lo que la civilización humana creó en los últimos 8 segundos de 530 million years''} & Esta frase usa la inteligencia animal para criticar la narrativa estrecha de AGI, y nos recuerda que no debemos confundir las tareas simbólicas de puntaje alto con la inteligencia misma. Sobrevivir, percibir y actuar en el mundo real es mucho más difícil que varios benchmarks abstractos. \\
\textit{``I am not the special one, I am the normal one''} & Esto es tanto un rechazo al mito individual como una afirmación del «espíritu de batería». La comunidad de investigación no necesita a un puñado de genios mitificados, sino a personas dispuestas a suministrar energía a largo plazo, tanto a los problemas como a los demás. \\
\textit{``No te preocupes por un point estimate... toda evaluación al final será una integral''} & Resume la actitud de Xie Saining hacia la revisión por pares, la aceptación de artículos, los premios y el éxito o fracaso a corto plazo. La evaluación parcial es un muestreo con mucho ruido; lo que verdaderamente importa es el impacto acumulado en un lapso de tiempo más largo. \\
\textit{``Cada persona es una variable en este mundo; es posible que tú seas justo la variable más importante''} & Este es uno de sus estímulos más directos para los investigadores jóvenes. Frente a un sistema altamente competitivo y con una fuerte dependencia de la trayectoria, cada persona debe conservar la voluntad de entrar activamente al campo, buscar activamente oportunidades y asumir activamente el riesgo. \\
\bottomrule
\end{longtable}

Si se ponen estas diez citas una junto a otra, se advierte que en realidad todas apuntan a la misma proposición central: los problemas que verdaderamente valen la pena dedicarles la vida entera no suelen estar donde la evaluación de corto plazo es más clara, sino en los lugares donde todavía es necesario que tú mismo entres, que tú mismo explores a tientas y que tú mismo asumas la incertidumbre. Ya sea que hable de visión, de método de investigación, de emprendimiento o de vida, Xie Saining nunca se aleja de esta proposición.

\end{document}
